% Options for packages loaded elsewhere
\PassOptionsToPackage{unicode}{hyperref}
\PassOptionsToPackage{hyphens}{url}
\documentclass[
  ignorenonframetext,
]{beamer}
\newif\ifbibliography
\usepackage{pgfpages}
\setbeamertemplate{caption}[numbered]
\setbeamertemplate{caption label separator}{: }
\setbeamercolor{caption name}{fg=normal text.fg}
\beamertemplatenavigationsymbolsempty
% remove section numbering
\setbeamertemplate{part page}{
  \centering
  \begin{beamercolorbox}[sep=16pt,center]{part title}
    \usebeamerfont{part title}\insertpart\par
  \end{beamercolorbox}
}
\setbeamertemplate{section page}{
  \centering
  \begin{beamercolorbox}[sep=12pt,center]{section title}
    \usebeamerfont{section title}\insertsection\par
  \end{beamercolorbox}
}
\setbeamertemplate{subsection page}{
  \centering
  \begin{beamercolorbox}[sep=8pt,center]{subsection title}
    \usebeamerfont{subsection title}\insertsubsection\par
  \end{beamercolorbox}
}
% Prevent slide breaks in the middle of a paragraph
\widowpenalties 1 10000
\raggedbottom
\AtBeginPart{
  \frame{\partpage}
}
\AtBeginSection{
  \ifbibliography
  \else
    \frame{\sectionpage}
  \fi
}
\AtBeginSubsection{
  \frame{\subsectionpage}
}
\usepackage{iftex}
\ifPDFTeX
  \usepackage[T1]{fontenc}
  \usepackage[utf8]{inputenc}
  \usepackage{textcomp} % provide euro and other symbols
\else % if luatex or xetex
  \usepackage{unicode-math} % this also loads fontspec
  \defaultfontfeatures{Scale=MatchLowercase}
  \defaultfontfeatures[\rmfamily]{Ligatures=TeX,Scale=1}
\fi
\usepackage{lmodern}
\ifPDFTeX\else
  % xetex/luatex font selection
\fi
% Use upquote if available, for straight quotes in verbatim environments
\IfFileExists{upquote.sty}{\usepackage{upquote}}{}
\IfFileExists{microtype.sty}{% use microtype if available
  \usepackage[]{microtype}
  \UseMicrotypeSet[protrusion]{basicmath} % disable protrusion for tt fonts
}{}
\makeatletter
\@ifundefined{KOMAClassName}{% if non-KOMA class
  \IfFileExists{parskip.sty}{%
    \usepackage{parskip}
  }{% else
    \setlength{\parindent}{0pt}
    \setlength{\parskip}{6pt plus 2pt minus 1pt}}
}{% if KOMA class
  \KOMAoptions{parskip=half}}
\makeatother
\setlength{\emergencystretch}{3em} % prevent overfull lines
\providecommand{\tightlist}{%
  \setlength{\itemsep}{0pt}\setlength{\parskip}{0pt}}
% Slide layout defaults for warning-clean scientific decks.
\usepackage{etoolbox}
\IfFileExists{xurl.sty}{\usepackage{xurl}}{}
\IfFileExists{seqsplit.sty}{\usepackage{seqsplit}}{\newcommand{\seqsplit}[1]{#1}}
\protected\def\breakseq#1{\seqsplit{#1}}
\protected\def\breaktt#1{\begingroup\ttfamily\seqsplit{#1}\endgroup}
\setlength{\emergencystretch}{6em}
\tolerance=5000
\hbadness=10000
\hfuzz=1pt
\setlength{\tabcolsep}{2pt}
\AtBeginEnvironment{longtable}{\tiny\renewcommand{\arraystretch}{0.86}\setlength{\tabcolsep}{1pt}}
\AtBeginEnvironment{tabular}{\tiny\renewcommand{\arraystretch}{0.86}\setlength{\tabcolsep}{1pt}}

% Natbib and cross-reference fallbacks — slides don't load natbib
% or cleveref, but combined-PDF manuscript prose may emit these
% commands. The fallback renders citations as a bracketed key list
% and cross-references as detokenized labels so slides don't fail on
% undefined control sequences or raw underscores. \providecommand is
% a no-op if packages load later.
\providecommand{\citep}[1]{[#1]}
\providecommand{\citet}[1]{#1}
\providecommand{\citealp}[1]{#1}
\providecommand{\citeauthor}[1]{#1}
\providecommand{\citeyear}[1]{#1}
\providecommand{\cref}[1]{\texttt{\detokenize{#1}}}
\providecommand{\Cref}[1]{\texttt{\detokenize{#1}}}

\newtheorem{proposition}{Proposition}
\newtheorem{hypothesis}{Hypothesis}
\usepackage{bookmark}
\IfFileExists{xurl.sty}{\usepackage{xurl}}{} % add URL line breaks if available
\urlstyle{same}
% fallback for those not using the hyperref driver hyperxmp:
\makeatletter
\@ifundefined{xmpquote}{\newcommand{\xmpquote}[1]{#1}}{}
\makeatother
\hypersetup{
  hidelinks,
  pdfcreator={LaTeX via pandoc}}

\author{\texorpdfstring{}{}}
\date{}

\begin{document}

\section{Abstract}\label{sec:abstract}

\begin{frame}[fragile,allowframebreaks]{Abstract}
This paper documents \breaktt{template\_search\_project}, the
literature-search exemplar shipped with the
\href{https://github.com/docxology/template}{Research Project Template}.
The project demonstrates \textbf{two configurable, reproducible
pipelines} sharing the same configuration file and the same
\breaktt{infrastructure/search/} + \breaktt{infrastructure/reference/}
modules. The standard pipeline
(\breaktt{scripts/run\_search\_pipeline.py}) handles a single
\texttt{SearchQuery} end-to-end. The deep-search pipeline
(\breaktt{scripts/run\_deep\_search.py}, see \breakseq{[@sec:deep\_search]})
fans out across a list of keywords (each capped at 100 papers per
keyword from \breaktt{deep\_search.max\_results\_per\_keyword} in
\breaktt{manuscript/config.yaml}), fully enriches every paper with its
abstract and PDF fulltext, and (optionally) uses the local LLM to write
a multi-section reading note for every paper. When a deep-search
aggregate exists, the latest run covered \textbf{3} keyword(s) with
\textbf{} unique paper(s) after cross-keyword deduplication. Both turn a
free-text topic into:

\begin{enumerate}
\tightlist
\item
  a deduplicated, year-filtered set of papers drawn from arXiv,
  Crossref, optional local corpora, and (opt-in)
  \href{https://paperclip.gxl.ai/}{Paperclip};
\item
  a Pandoc-compatible \texttt{references.bib} byte-identical in style to
  the canonical exemplar in
  \href{../../template_code_project/manuscript/references.bib}{\breaktt{template\_code\_project}}
  (file \breaktt{manuscript/references.bib});
\item
  cached abstracts and (optionally) extracted PDF full text, written to
  disk under stable per-paper identifiers; and
\item
  an LLM-synthesised reading report assembled from per-paper analyses
  and a cross-corpus thematic synthesis, all produced by a local Ollama
  model with pinned seed and temperature.
\end{enumerate}

All discovery logic lives in \breaktt{infrastructure/search/literature/}
(\href{https://github.com/docxology/template/tree/main/infrastructure/search/literature}{source
on GitHub}); all export logic lives in
\breaktt{infrastructure/reference/citation/}
(\href{https://github.com/docxology/template/tree/main/infrastructure/reference/citation}{source
on GitHub}); LLM synthesis reuses the existing
\breaktt{infrastructure/llm/}
(\href{https://github.com/docxology/template/tree/main/infrastructure/llm}{source
on GitHub}) bridge. The project itself contains only thin orchestration,
manuscript prose, and a test suite --- perfectly mirroring the
\textbf{two-layer architecture} the template enforces.

The motivating concern is \emph{reproducibility}: a query at time
\(t_0\) should produce the same results at time \(t_1\) unless the cache
is explicitly invalidated. This is achieved by deterministic search
caching keyed on canonical query identity, on-disk caching of every
fetched abstract / PDF, and pinned LLM seeds. The same
\breaktt{manuscript/config.yaml} that drives the pipeline is also the
only configuration any reviewer needs.

\textbf{Run snapshot.} With the bundled \breaktt{manuscript/config.yaml},
the most recent pipeline execution evaluated the query
\emph{``reproducible research optimization''} against local, returned 6
deduplicated paper(s) (4 carrying a DOI, 6 carrying an abstract), and
recorded backend errors: none. Resolve \texttt{\{\{…\}\}} tokens by
running \breaktt{scripts/z\_generate\_manuscript\_variables.py} after
\breaktt{run\_search\_pipeline.py}; the script writes
\breaktt{output/data/manuscript\_variables.json} and resolved markdown
under \breaktt{output/manuscript/}, which the PDF-rendering stage prefers
when present.

\textbf{Keywords:} literature search, BibTeX automation, reproducible
research, local LLM synthesis, scientific infrastructure
\end{frame}

\end{document}
